\preludeandfugue{(No) Lingua Franca}
\contributor{Irena Peterson Pavliuk}

\settowidth{\versewidth}{My father talks to me about politics from the clouds,}
\begin{verse}[\versewidth]
  
% (No) Lingua Franca\\!
  
% \textbf{1}\\!
  
Oh, leave us without your curiosity.\\
The colonial tongue is no lingua franca for us.\\!

We are not even the children of the murdered ones.\\
We are the children of those who kept silent\\
and, thus, survived.\\!

My mother lives in deep countryside,\\
from where her voice is not heard.\\
My father talks to me about politics from the clouds,\\
he did not live to see the revolution.\\!

I am the child of those who kept silent.\\
Chances are you are as well.\\
Is this why you are all so quiet,\\
and I barely know you,\\
but also have known you for lifetimes?\\!

\[***\]
%\textbf{2}\\!
\clearpage

Songs of the forests\\
\qquad we share,\\
flapping of the wings over the fields\\
\qquad we share,\\
dark seas where we feel at home but with so much pain---\\
because what is home and for how long?\\!

We share\\
\qquad the habit of being quiet,\\
we share\\
\qquad the custom of saving bread for tomorrow.\\!

We share\\
\qquad being seen as those behind the wall,\\
we share\\
\qquad being not seen at all.\\!

\vspace*{1cm}
%\[-\]
%\textbf{3}\\!

Emergency backpacks\\
in the cellars \\
of our homes---\\
we're not burdened with the material things---\\
it's called minimalism;\\
it's now in fashion.\\!

We play games in the woods;\\
we jump over fires on Midsummer's Eves;\\
we jump over fires---always!\\!

We played games in the woods:\\
I learnt to make a perfect dugout when I was five---\\
my grandpa didn’t like it.\\!

\clearpage
%\[-\]
%\textbf{4}\\!

The remnants of self-preserving:\\
they read as nationalism.\\
I read it as ``try to shame me''.\\!

Oh, take me to my Midsummer's fire and try to shame me.\\!

Oh, throw me into the dark deep water---it knows us so well.\\!

The sweetest song my grandmother sang.\\
The sweetest song your grandmother sang.\\!

The water does not speak the colonial tongue,\\
but it knows our songs too well\\
and needs no lingua franca.\\!

\end{verse}
